Algorithmic redistricting presents a paradox: a process that receives no information
about incumbent legislators nonetheless produces outcomes that profoundly affect their
electoral prospects. This paper synthesises the findings of I.1--I.4 into a unified
framework for understanding how the \textsc{bisect} recursive-bisection algorithm
interacts with congressional incumbency across three states---North Carolina ($k=14$),
Wisconsin ($k=8$), and Texas ($k=38$)---under 2020 Census data. We identify three
distinct mechanisms through which a politically blind algorithm affects incumbents:
geographic restructuring that pairs incumbents in the same district, changes in
partisan composition that alter the safe/competitive balance, and boundary shifts
that reassign constituent bases. We introduce the incumbency-neutral baseline as the
expected value of each incumbency metric under random redistricting given the state's
geographic constraints, and show that \textsc{bisect} outcomes fall within this
baseline range on all three metrics across all three states. Enacted 2022 congressional
maps, by contrast, produce significantly lower incumbent pairing rates and higher
safe-seat fractions than the algorithmic baseline in all three states, consistent with
deliberate use of incumbent address data that the \textsc{bisect} algorithm does not
receive. The legal framework governing algorithmic treatment of incumbency is asymmetric:
courts have held that incumbency protection is a permissible redistricting criterion,
but no court has required it, and algorithmic indifference to incumbency is
constitutionally sound. These findings establish \textsc{bisect} as an incumbency-neutral
baseline suitable for expert witness testimony in redistricting litigation.
